% Source presentation for register row N-c-028.

\subsection{McMillan's interpolation-based finite-state safety method}
\label{sec:p09-presentation-n-c-028}

For an inconsistent pair \((A,B)\), McMillan uses a Craig interpolant \(P\)
satisfying
\[
 A\models P,\qquad P\land B\models\bot,\qquad
 \operatorname{Var}(P)\subseteq
 \operatorname{Var}(A)\cap\operatorname{Var}(B).
\]
The source constructs \(P\) from a resolution refutation by labelling an
\(A\)-root with the disjunction of its shared literals, labelling a \(B\)-root
with \(\top\), and combining predecessor labels by disjunction when the pivot
is local to \(A\) and by conjunction otherwise
\cite[def.~2 and thm.~2]{McMillan2003}.

For the clauses
\[
 A=\{p\},\qquad
 B=\{\neg p\lor q,\ \neg q\},
\]
resolve the two \(B\)-clauses on \(q\) to obtain \(\neg p\), then resolve with
the \(A\)-clause to obtain the empty clause.  The formula \(P=p\) is an
interpolant: \(A\models p\), the set \(B\cup\{p\}\) is inconsistent, and \(p\)
is the only variable common to the two sides.  This tiny proof shows why a
refutation-producing SAT procedure can return a state formula rather than only
the word UNSAT.  It also calculates McMillan's labelling rule exactly: the
\(A\)-root receives label \(p\), both \(B\)-roots receive \(\top\), resolving
the \(B\)-roots on the \(B\)-local variable \(q\) gives
\(\top\land\top=\top\), and the final resolution on the shared variable \(p\)
gives \(p\land\top=p\).

The transition-system presentation is
\[
 M=(I(V),T(V,V'),F(V)),
\]
with total Boolean transition relation \(T\).  A definitional CNF may introduce
fresh variables \(U\), but it must preserve projection:
\[
 \exists U\,\operatorname{Cnf}(f,U)\ \leftrightarrow\ f.
\]
For a current reachability overapproximation \(R\) and a fixed \(k>0\), one
interpolation iteration uses disjoint CNF auxiliary variables and the partition
\[
\begin{aligned}
 A={}&\operatorname{Cnf}
   \bigl(R(W_{-1})\land T(W_{-1},W_0),U_1\bigr),\\
 B={}&\operatorname{Cnf}
   \left(
    \bigwedge_{0\leq i<k}T(W_i,W_{i+1})
    \land\bigvee_{0\leq i\leq k}F(W_i),U_2
   \right).
\end{aligned}
\]
The two sides share only the boundary-state variables \(W_0\).  The
\(A\)-side describes a successor of \(R\); the \(B\)-side describes a suffix of
the fixed positive horizon containing a bad state.  Because \(T\) is total, a
path that first reaches a bad state in at most \(k\) steps can be extended
through the unused state copies, and every satisfying assignment to \(B\)
contains such a prefix.  An interpolant \(P(W_0)\) therefore overapproximates
the one-step successors of \(R\) while excluding every boundary state that can
reach a bad state in at most \(k\) steps.

Take the one-bit system
\[
 I(x)=\neg x,\qquad
 T(x,x')=(x'\leftrightarrow x),\qquad
 F(x)=x.
\]
For horizon \(k=1\) and current approximation \(R=I\), the first side entails
\(\neg x_0\).  The suffix side requires
\[
 (x_1\leftrightarrow x_0)\land(x_0\lor x_1),
\]
which is inconsistent with \(\neg x_0\).  Thus \(P(x_0)=\neg x_0\) is an
interpolant.  After renaming \(x_0\) to \(x\), the update
\(R:=R\lor P\) changes nothing.  The fixed point proves the inductive invariant
\(\neg x\), which contains every initial state, is closed under \(T\), and
excludes \(F\).

The invariant behind that calculation is general.  Initially \(R=I\).  From
\(A\models P\), every one-step successor of an \(R\)-state satisfies \(P\), so
the update \(R':=R\lor P\) still overapproximates all states reached so far.
From \(P\land B\models\bot\), no \(P\)-state reaches \(F\) within the protected
suffix; in particular \(P\land F\) is inconsistent.  Hence every update
preserves \(R'\land F\models\bot\).  If the fixed-point check
\(P\models R\) succeeds, then
\[
 R(V)\land T(V,V')\models P(V')\models R(V'),
\]
so the old \(R\) is an inductive invariant containing \(I\) and excluding
\(F\).

Three correctness cases must remain distinct
\cite[thms.~3--4]{McMillan2003}.

\begin{enumerate}
\item A satisfying bounded formula on the first iteration supplies a real bad
run from an initial state.
\item At a later iteration, \(R\) may include unreachable states, so a
satisfying suffix may be spurious; the procedure abandons that horizon rather
than reporting a counterexample.
\item An interpolation fixed point is an inductive invariant and therefore an
unbounded safety proof.
\end{enumerate}

The introductory prose permits a generic \(k\geq0\), but the displayed
\textsc{FiniteRun} procedure, Theorem 3, and the termination development use
\(k>0\).  Definition~\ref{def:p09-interpolant-model-checking} retains the proved
positive hypothesis.  Theorem 4 proves termination by choosing the reverse
depth as the horizon.  To retain the procedure's positive-horizon hypothesis
also when that depth is \(0\),
Theorem~\ref{thm:p09-interpolant-model-checking} uses the edge-safe choice
\(k=\max\{1,\operatorname{rd}(M,F)\}\).  Indeed, after the first unsatisfiable
iteration the new approximation contains no state that reaches \(F\) within
\(k\) steps.  When \(k\) is at least the reverse depth, any state that reaches
\(F\) at all has a shortest such path of length at most \(k\); the new
approximation therefore cannot reach \(F\) at any length.  Every later SAT
query is consequently unsatisfiable.  Unless the fixed-point implication has
already succeeded, \(R\lor P\) denotes a strict enlargement of \(R\).  The
Boolean state space is finite, so only finitely many strict enlargements are
possible and the procedure reaches a fixed point.

The source theorem is about finite Boolean safety/reachability and
resolution-derived interpolants; it is not a general temporal-logic theorem or
a completeness claim about arbitrary interpolant choices and SAT
implementations.
